\documentclass[12pt]{article}
\usepackage[utf8]{inputenc}
\usepackage[T1]{fontenc}
\usepackage{hyperref}
\title{Washo Language Resources for K-12 Curriculum Development}
\author{Compiled by Claude Sonnet 5 (claude-sonnet-5) for Andr\'as Kornai}
\date{}
\begin{document}
\maketitle

\section{Introduction}
Washo (ISO 639-3 \texttt{was}, Glottocode wash1253) is a language isolate of the Washoe Tribe of Nevada and California. It is severely endangered. Sources found in search give figures from about eight to twenty fluent speakers, mostly elderly. This report is short because little was found in search.

\section{Primary Educational Resources}
No classroom textbooks or primers were found. Sources report that an immersion school opened in 1997 and later closed, and that the language programme now operates through the tribe's Cultural Resource Department:
\begin{itemize}
\item American Indian Magazine, ``Where Wasiw is Spoken'': \url{https://www.americanindianmagazine.org/story/where-wasiw-spoken}
\end{itemize}

\section{Linguistic and Reference Materials}
\begin{itemize}
\item Wikipedia overview of the Washo language (background only): \url{https://en.wikipedia.org/wiki/Washo_language}
\item An eScholarship item on Washo, found in search; I did not review its contents: \url{https://www.escholarship.org/uc/item/4zd060kz}
\item A 2007 news report on elders helping compile an online dictionary with a linguist. I could not confirm whether the dictionary was finished or where it is hosted: \url{https://indianz.com/News/2007/005030.asp}
\end{itemize}

\section{Religious and Community Resources}
The Washoe Tribe's language programme is run through its Cultural Resource Department. The search did not return a tribal web page for it, so no URL is given here beyond the articles above.

\section{Digital and Audio Resources}
No verified audio or digital resources were found in this search.

\section{Curriculum Development Context}
Washo is not known to be taught in schools at present. A curriculum would need to be built in direct partnership with the Washoe Tribe and its remaining speakers. A first step is to find out what happened to the online dictionary project and what the tribe's programme currently offers.

\end{document}